\subsection{Payout Annuities}


\content{
        Formulas so far: \\
        \fbox{\begin{minipage}{\textwidth}
        For these formulas, we have the following: 
        $P_0$ is the pricipal. 
        $r$ is the interest rate. 
        $t$ is the number of time intervals (specified by problem) which
        have passed. 
        $A$ is the amount paid back.
        $I$ is the interest.

        \begin{center}
        \begin{tabular}{|c|c|}
        \hline
        \T\B Simple Interest &  $ A = P_0(1+r) $ \\
        \hline
        \T\B Simple Interest over Time & $  A = P_0(1+rt) $ \\
        \hline
        \end{tabular}
        \end{center}
        \end{minipage}}
        \fbox{\begin{minipage}{\textwidth}
        For these formulas, we have the following: 
        $P_0$ is the pricipal. 
        $r$ is the annual interest rate.
        $t$ is the number of years which have passed. 
        $A$ is the total amount.
        $k$ is the number of times interest is compounded per year. 

        \begin{center}
        \begin{tabular}{|c|c|}
        \hline
        \T\B Compound Interest & $ A =
        P_0\left(1+\frac{r}{k}\right)^{kt} $ \\
        \hline
        \T\B Savings Annuity & $ A =
        d\frac{\left(\left(1+\frac{r}{k}\right)^{kt}-1\right)}{\frac{r}{k}} $\\
        \hline
        \end{tabular}
        \end{center}
        \end{minipage}}
}{% presentation
}{% nopresentation
}{% comments
}

\content{
        \begin{definition} (Payout Annuities) A payout annuity is when you have
        some amount of money in an account, and you make regular deposits from
        it, leaving the rest to earn interest. The formula is as follows.
        $$ P_0 =
        d\frac{\left(1-\left(1+\frac{r}{k}\right)^{-kt}\right)}{\frac{r}{k}} $$
        \end{definition}
}{% presentation
}{% nopresentation
}{% comments
}

\content{
        \begin{example} After retiring, you want to be able to take \$1000 every month for a total of 20 years from
        your retirement account. The account earns 6\% interest. How much will you need in your
        account when you retire? How much will you withdraw from the account
        total? How much of that money is interest? 
        \end{example}
}{% presentation
\vspace{3in}
}{% nopresentation
\vspace{3in}
}{% comments
}

\content{
        \begin{example} (Negative Exponents on a Calculator) Depending on your
        calculator, the steps which are needed to calculate something like
        $1.005^{-240}$ may differ.
        \end{example}
}{% presentation
\vspace{3in}
}{% nopresentation
\vspace{3in}
}{% comments
}

\content{
        \begin{example} You know you will have \$500,000 in your account when you retire. You want to be able to
        take monthly withdrawals from the account for a total of 30 years. Your retirement account
        earns 8\% interest. How much will you be able to withdraw each month?
        \end{example}
}{% presentation
\vspace{3in}
}{% nopresentation
\vspace{3in}
}{% comments
}

\content{
        \begin{example} A donor gives \$100,000 to a university, and specifies that it is to be used to give annual
        scholarships for the next 20 years. If the university can earn 4\% interest, how much can they
        give in scholarships each year?
        \end{example}
}{% presentation
\vspace{3in}
}{% nopresentation
\vspace{3in}
}{% comments
}

\content{
        \begin{example} How much money will I need to have at retirement so I
        can withdraw \$60,000 a year for 20 years from an account earning 8\% 
        compounded annually? How much money will be withdrawn from the account?
        How much of that money is interest? 
        \end{example}
}{% presentation
\vspace{4in}
}{% nopresentation
\vspace{3in}
}{% comments
}
